\section{Review of Existing Alternatives}

This section discusses how farmers in Sinacaban currently cope with the lack of a centralized agricultural information system. It also reviews existing digital platforms related to farming, produce trading, and agricultural coordination, and explains how the proposed system addresses gaps that remain in these alternatives.

Several digital initiatives and information systems have been developed to improve market linkages and agricultural coordination among farmers in the Philippines and in other developing countries. One notable example is the Department of Agriculture's e-Kadiwa program, which aims to connect local producers directly with consumers through an online marketplace. Similarly, the DA eMarketplace provides a national digital platform for trading agricultural products. These systems show the government's efforts to support digital transformation in the agricultural sector. However, studies have shown that the use of these platforms in rural communities remains limited because of poor internet connectivity, low digital literacy, and the lack of localized features for smallholder farmers \citep{Briones2023,Pastolero2022}.

In practice, many smallholder farmers continue to rely on social media platforms such as Facebook Marketplace and community group chats to sell their products. These tools are accessible and familiar to many users, but they lack structured features for managing seasonal harvests, recording crop data, monitoring produce availability, and supporting institutional coordination. Transactions through these informal channels are often unmonitored, making it difficult for farmers and agriculture personnel to track prices, quantities, and the flow of goods. Similar observations were reported by \citet{Ma2024}, who noted that unstructured digital exchanges may limit long-term coordination and reduce the participation of institutional actors. This situation shows the need for dedicated agricultural information systems that combine accessibility with specialized features for coordination, transparency, and market linkage support \citep{JPAL2023,Qiang2012}.

Outside the Philippines, agricultural platforms such as FarmLogs, AgriDigital, and other commercial farm-management applications have been developed to address similar concerns. These tools commonly focus on record-keeping, logistics management, inventory monitoring, and price transparency. While they can improve efficiency and decision-making, many of them are designed for commercial-scale farms or areas with stable internet access. As a result, they may not be suitable for rural barangays in the Philippines, where most farmers operate on a small scale, have limited financial resources, and experience intermittent connectivity \citep{Ma2024}.

Previous studies emphasize that information asymmetry and weak digital connectivity remain major barriers to the productivity and profitability of rural farmers \citep{Briones2023}. \citet{Mopera2016} reported that post-harvest losses in the Philippines can reach up to 42\% for vegetables and 28\% for fruits, mainly due to inadequate logistics and poor coordination among market actors. The World Bank and \citet{Qiang2012} also noted that digital technologies can improve agricultural performance when they are designed to be inclusive, accessible, and appropriate to the local context \citep{WorldBank2013}. However, many existing systems focus mainly on market transactions and do not fully address the seasonal nature of farming, where timing, planning, and local monitoring are also important.

The \textit{Digital Farmers' Produce System} addresses these gaps by providing a localized, cloud-based platform designed for farmers, buyers, and the Municipal Agriculture Office of Sinacaban. Unlike national platforms and commercial applications, the system focuses on municipal-level agricultural coordination. It allows farmers to post available produce, indicate quantity and harvest information, send resource requests, and communicate with potential buyers. Buyers can view available crops in the area and contact farmers directly, while agriculture personnel can monitor crop listings, resource requests, and basic stock information.

The system introduces a three-way coordination model that connects farmers, buyers, and the Municipal Agriculture Office through one web-based platform. By combining produce visibility, resource request management, buyer communication, and basic monitoring tools, the system provides a practical and localized alternative to existing platforms. Although it does not include online payment, delivery tracking, or full logistics management, it addresses the need for a simple, user-friendly, and locally focused agricultural information system for rural municipalities.

\begin{table}[ht]
\centering
\caption{Summary of Existing Alternatives}
\label{table:existing-alternatives-summary}
\renewcommand{\arraystretch}{1.3}
\setlength{\tabcolsep}{4pt}
\footnotesize
\begin{tabular}{|>{\raggedright\arraybackslash}p{2.8cm}|
                    >{\raggedright\arraybackslash}p{1.8cm}|
                    >{\raggedright\arraybackslash}p{3.0cm}|
                    >{\raggedright\arraybackslash}p{3.0cm}|
                    >{\raggedright\arraybackslash}p{3.0cm}|}
\hline
\textbf{System (License)} &
\textbf{Platform} &
\textbf{Brief Description} &
\textbf{Basic Functionalities} &
\textbf{Relation to the Proposed System} \\
\hline\hline

\textbf{Proposed System: Digital Farmers' Produce System} &
Web &
Cloud-based agricultural platform for farmers, buyers, and the Municipal Agriculture Office of Sinacaban &
Produce posting, buyer viewing, messaging, resource requests, dashboard summaries, and role-based access &
Designed specifically for local agricultural coordination in Sinacaban \\
\hline

DA e-Kadiwa (Government platform) &
Web &
National online marketplace connecting agricultural producers and consumers &
Online product catalog, ordering, and basic market linkage support &
Useful for broader market access but not focused on municipal-level crop monitoring and resource requests \\
\hline

DA eMarketplace (Government platform) &
Web &
Digital platform for trading agricultural goods &
Product posting, price display, and online agricultural trading support &
Similar in supporting produce visibility, but less focused on localized farmer monitoring and LGU coordination \\
\hline

Facebook Marketplace / Group Chats (Proprietary platform) &
Web and Mobile App &
Informal selling and messaging channels commonly used by farmers and buyers &
Product posting, chat-based negotiation, and basic location-based search &
Accessible and familiar, but lacks structured crop records, request management, and institutional monitoring \\
\hline

Commercial farm-management applications such as FarmLogs and AgriDigital (Commercial / Subscription) &
Web and Mobile App &
Farm-management tools mainly designed for larger farms and commercial agricultural operations &
Field records, inventory tracking, logistics management, and farm analytics &
Feature-rich but may be costly, complex, and not localized for smallholder farmers in rural municipalities \\
\hline
\end{tabular}
\end{table}